\section{Reflection positivity of the original limit}\label{sec:completion}

\subsection{Convex separation and the nine mirror families}\label{subsec:diag-separation}

Fix an integer $c_0$. Let $\mathcal R(c_0)$ be the set of
laws on $\Omega$ satisfying Gaussian domination and both
invariance and reflection positivity at
$(x_1-x_2,c_0)$. This is convex and weakly compact. Gaussian
domination gives a common tail
$\lambda(|F_\sigma|\ge R)\le2e^{-R^2/2}$ and hence uniform
tightness. It is weakly closed by truncation of its exponential
integrals. Mirror invariance and positivity against continuous
local tests are also closed, and the approximation in \cref{subsec:coordinate-limits}
extends the latter to the full open-side $L^2$ space. Each condition is affine
or an affine inequality in the law, proving convexity.

We claim that, for every real local trigonometric polynomial
$O$ and every cubic limit $\nu$ along an arbitrary sequence of sides tending to infinity, some $\nu_0\in\mathcal R(c_0)$ satisfies
\eqref{eq:diag-dominating-state}. The polynomial need not
initially fit in a cell having its lower wall at $c_0$.
Let $\supp O\subset y_0+[-r,r]^3$, with $r\ge1$, and choose
an admissible even width $h>8r+8$. Choose an even integer
$a$ in
\[
 (c_0+2r-\ell_-(y_0),\ c_0+h-2r-\ell_-(y_0)).
\]
The interval has length greater than two. For $v=ae_1$ and
$y=y_0+v$ its defining inequalities place the translated
box between the two $\ell_-$ walls. Choose integer $c_+$
and half-integer $c_3$ so that
\[
 c_+<\ell_+(y)-2r<\ell_+(y)+2r<c_++h,
 \qquad c_3<y_3-r<y_3+r<c_3+h.
\]
The available intervals have length greater than one.
The box therefore lies in the open diagonal cell, since
$|\ell_\pm(x)-\ell_\pm(y)|\le2r$ on it. Form the tilted
law with the translated polynomial $O^{\tau_v}$.

Choose companion sides by \cref{lem:companions-exist} and form a
comparison state on those sides, subordinate to the original volume
subsequence defining $\nu$. It belongs to $\mathcal R(c_0)$,
and \cref{thm:comparison-any-side} gives
$\nu(O^{\tau_v})\le\nu_0(O^{\tau_v})$. The original cubic
laws are invariant under every lattice translation, and
this invariance passes to $\nu$. Equation
\eqref{eq:diag-even-translation} gives the same invariance
under this particular $v$ for $\nu_0$. Hence
\begin{equation}\label{eq:diag-dominating-state}
 \nu(O)=\nu(O^{\tau_v})
       \le\nu_0(O^{\tau_v})=\nu_0(O).
\end{equation}
The comparison state may depend on $O$. Convex separation
is what turns these individual inequalities into reflection
positivity of the original law.

\begin{lemma}\label{lem:diag-separation}
Let $\mathcal R$ be a convex weakly compact set of laws on
$\Omega$. If for every real local trigonometric polynomial
$O$ some $\lambda\in\mathcal R$ satisfies
$\nu(O)\le\lambda(O)$, then $\nu\in\mathcal R$.
\end{lemma}

\begin{proof}
Give the real vector space of finite signed measures on
$\Omega$ the weak topology defined by integration against
these polynomials. It is locally convex and Hausdorff:
if every such integral of a signed measure vanishes, its
finite-coordinate marginals have zero Fourier transform,
so its positive and negative parts agree on cylinder sets
and hence on every Borel set. Its continuous linear
functionals are precisely integration against the test
space; a continuous linear functional for a topology
defined by these seminorms depends on finitely many of
them, and their finite real span is still that space.
Since the tests are bounded and continuous, the inclusion
from the usual weak topology of laws is continuous.
Thus $\mathcal R$ is a compact convex subset of this
locally convex Hausdorff vector space. It is nonempty
by the hypothesis at $O=0$.

If $\nu$ were outside $\mathcal R$, strict separation
\cite[Theorem~6.2.9]{Peterson2017} of
a compact convex set and a disjoint closed point would
give a continuous linear functional with
$\max_{\lambda\in\mathcal R}\lambda(O)<\nu(O)$ for some
test $O$. Its maximum exists by compactness. This contradicts
the assumed inequality for that $O$.
\end{proof}

\begin{proof}[Proof of \Cref{thm:main}]

The coordinate part of the theorem, including tightness, Gaussian
domination, all lattice symmetries and coordinate reflection positivity,
was proved in \cref{subsec:coordinate-limits} for arbitrary integer
sides. Let $\nu$ be any limit along a prescribed original sequence.
For the diagonal mirror $(x_1-x_2,c_0)$,
\eqref{eq:diag-dominating-state} holds for every real
trigonometric polynomial. \Cref{lem:diag-separation}
therefore gives $\nu\in\mathcal R(c_0)$, proving its
diagonal positivity. This is valid for every integer
$c_0$.

A signed permutation transports this diagonal mirror
and its positive half-space to any plane
$x_\mu\pm x_\nu=c$. The corresponding reflection is the
conjugate of the original one. Pull a positive-side test
back by this isometry; invariance of $\nu$ identifies
its reflected form with the already nonnegative one.
This gives the six diagonal families and completes
the proof.
\end{proof}

